\begin{figure}[!htbp]
    \centering
    \vspace{-0.2cm}
    \includegraphics[width=\linewidth]{figures/ubtree_overview_demo.pdf}
    \vspace{-0.2cm}
    \caption{\textbf{(a) Speedups on Qwen3-8B:} Speedup over autoregressive decoding across seven benchmarks on Qwen3-8B instruct at $T_{\mathrm{infer}}=1$.
     \textbf{(b) Speedups on Ling3-Flash-124B:} Serving speedup in SGLang on Ling3-Flash-124B at $T_{\mathrm{infer}}=1$.
    \textbf{(c) Long-context acceptance:} Acceptance length $\tau$ on Ling3-Flash-124B for LongBench-v2 and SWE-bench over different context lengths. UBTree outperforms competing methods across most benchmarks and model scales.}
    \label{fig:ubtree-teaser}
\end{figure}

\section{Introduction}
Speculative decoding~\citep{leviathan2023fast,chen2023accelerating} offers a lossless approach to accelerating large language models (LLMs): a lightweight drafter proposes multiple future tokens, and the target model verifies them together while preserving its output distribution. The benefit of speculative decoding depends not only on the acceptance length of the draft by the target model but also on the computational cost. Longer accepted drafts do not necessarily yield greater speedup if producing them incurs substantial latency.

Recent advances in speculative decoding have been driven by breakthroughs in parallel drafting~\citep{liu2026dart,chen2026dflash}. While their autoregressive counterparts~\citep{li2024eagle,li2026eagle} condition each draft token sequentially on preceding choices, parallel drafters predict all future positions in a drafting block with a single proposer-backbone forward pass. Parallelism reduces the sequential overhead of using larger proposer backbones which increases the acceptance length. However, parallel drafters lack draft diversity. Consequently, their performance degrades when high-entropy target distributions inherently produce different paths. Moreover, the missing token-wise dependency may worsen the degradation by mixing these different paths. Existing approaches try to complement this drawback by restoring token-wise dependencies to improve precisions of the single-path draft through lightweight autoregressive correction heads~\citep{huang2026domino,cheng2026dspark} or intra-backbone designs~\citep{inco2026dflash2}. However, these strategies do not fundamentally solve the diversity bottleneck because there is still only one single draft path covered. This particularly limits the further improvement of parallel drafters and their applications under scenarios with high entropy target distributions~\citep{devic2025trace,yang2026longspec}. \textbf{Draft diversity thus becomes a central bottleneck for current parallel speculative decoding.}

A natural solution to this bottleneck is tree drafting~\citep{miao2024specinfer,cai2024medusa,li2024eagle2}, which retains multiple draft paths as tree branches, increasing the chance that target verification matches one of them. By providing diverse yet plausible candidate drafts, tree drafting is more robust in high-entropy scenarios. Combining tree drafting and parallel drafters by constructing budgeted trees from pretrained block-parallel drafters achieves considerable speedups~\citep{ringel2026accelerating,li2026dartree}. Nevertheless, these methods exhibit two limitations. First, the tension between token-wise dependency and full parallelism constrains the tradeoff between draft quality and drafting latency. Existing tree-based methods use autoregressive heads to capture token-wise dependencies, which precludes full parallelization. Conversely, omitting token-wise dependency limits the acceptance length and, consequently, the end-to-end speedup. Second, tree-based parallel methods suffer from a \textit{train--inference mismatch}: they directly apply parallel drafters trained for single-path drafting to multi-path tree verification. Such drafters are optimized to produce the top draft for single-path verification, but are suboptimal for proposing multiple plausible paths for tree verification. Consequently, although lower-ranked alternatives in the draft tree can populate additional branches, they are not sufficiently optimized to align with the target model's alternatives or worth allocating verification budget to.

We address these limitations with \textbf{UBTree}, a tree-based drafter with a parallel bigram selector to model token-wise dependencies and tree-native training to close the train--inference gap. We incorporate two designs in UBTree:
\begin{itemize}
    \item \textbf{Architecture.} UBTree separates drafting into three complementary stages and keeps all neural computation parallel. First, we use DFlash~\citep{chen2026dflash} as a \textit{unigram proposer} to independently generate candidate tokens for every position in a drafting block. Second, a redesigned \textit{bigram selector} scores adjacent candidate-token pairs conditioned on the proposer hidden states. It derives predecessor and successor representations from frozen target embeddings through separate MLPs and learns depth-dependent weights for combining proposer logits and selector scores. Third, we use the resulting transition scores to construct the verification tree. By simplifying the token-wise dependency into bigram correlation, we overcome the challenge of parallelizing token-wise dependency in tree drafting.
    \item \textbf{Training.} Unlike other drafters trained on target-regenerated rollouts at the inference temperature, we use a renormalized KL training objective on high-temperature regenerated target trajectories to diversify the positive training labels. This \textit{tree-native} training strategy allocates more supervision signals to plausible alternative paths in the target distribution and distinguishes them better from low-quality alternative paths in tree construction, thereby relieving the train--inference mismatch in tree drafting.
\end{itemize}
We extensively evaluate UBTree on both academic and production-scale benchmarks. UBTree outperforms all existing methods on both academic-scale and production benchmarks. It achieves an average decoding speedup of $5.84$--$6.94\times$ over autoregressive decoding on Qwen3-4B and Qwen3-8B with increased acceptance lengths. It also rivals frontier baselines such as DSpark~\citep{cheng2026dspark} in the acceleration of production LLMs like Ling3-Flash-124B. Additionally, UBTree shows outstanding durability on long-context tasks compared to baselines. In summary, UBTree provides a frontier solution for accelerating LLMs in both academic settings and production environments.
